\frfilename{mt122.tex}
\versiondate{4.1.04}
\copyrightdate{1994}

\def\chaptername{Integrasi}
\def\sectionname{Definisi integral}

\newsection{122}

Saya menguraikan definisi integrasi biasa untuk fungsi bernilai real
yang didefinisikan pada suatu ruang ukur sebarang, beserta sifat-sifat
paling dasarnya.

\leader{122A}{Definisi}  Misalkan $(X,\Sigma,\mu)$ suatu ruang ukur.

\header{122Aa}{\bf (a)} Untuk setiap himpunan $A\subseteq X$, saya menuliskan
$\chi A$ untuk
{\bf fungsi karakteristik} dari $A$, yaitu fungsi
dari $X$ ke $\{0,1\}$ yang ditentukan dengan menetapkan $\chi A(x)=1$ jika $x\in A$, dan $0$
jika
$x\in X\setminus A$.   \cmmnt{(Tentu saja, notasi ini bergantung pada
dipahaminya himpunan `semesta' $X$ mana yang sedang ditinjau;
barangkali saya seharusnya menyebutnya `fungsi karakteristik dari $A$ sebagai
subhimpunan dari $X$'.)}
\cmmnt{Perhatikan bahwa} $\chi A$ bersifat $\Sigma$-terukur\cmmnt{, dalam
pengertian
121C di atas,} jika dan hanya jika $A\in\Sigma$\cmmnt{ (karena
$A=\{x:\chi A(x)>0\}$)}.

\header{122Ab}{\bf (b)} Sekarang, sebuah fungsi {\bf sederhana} pada $X$ adalah
fungsi
berbentuk
$\sum_{i=0}^n a_i\chi E_i$, dengan $E_0,\ldots,E_n$ merupakan himpunan
terukur yang mempunyai ukuran berhingga dan $a_0,\ldots, a_n$ termasuk dalam
$\Bbb R$.   \cmmnt{{\bf Peringatan!} Sebagian penulis memperbolehkan sembarang
himpunan $E_i$, sehingga sebuah fungsi sederhana pada $X$ adalah sembarang fungsi yang hanya
mengambil berhingga banyak nilai.}

\leader{122B}{Lema} Misalkan $(X,\Sigma,\mu)$ suatu ruang ukur.

(a) Setiap fungsi sederhana pada $X$ bersifat terukur.

(b) Jika $f$, $g:X\to\Bbb R$ adalah fungsi sederhana, maka $f+g$ juga
merupakan fungsi sederhana.

(c) Jika $f:X\to \Bbb R$ adalah fungsi sederhana dan $c\in\Bbb R$, maka
$cf:X\to\Bbb R$ merupakan fungsi sederhana.

(d) Fungsi nol konstan merupakan fungsi sederhana.

\proof{ (a) mengikuti fakta bahwa $\chi E$
terukur untuk setiap $E$ terukur, dan bahwa jumlah serta kelipatan skalar dari
fungsi-fungsi terukur bersifat terukur (121Eb-121Ec).  (b)-(d)
jelas.
}%end of proof of 122B

\vleader{42pt}{122C}{Lema} Misalkan $(X,\Sigma,\mu)$ suatu ruang ukur.

(a) Jika $E_0,\ldots,E_n$ merupakan himpunan terukur yang mempunyai ukuran berhingga,
terdapat himpunan terukur saling lepas $G_0,\ldots,G_m$ yang mempunyai ukuran berhingga
sedemikian sehingga setiap $E_i$ dapat dinyatakan sebagai gabungan dari beberapa $G_j$.

(b) Jika $f:X\to\Bbb R$ adalah fungsi sederhana, fungsi tersebut dapat dinyatakan dalam
bentuk $\sum_{j=0}^mb_j\chi G_j$, dengan $G_0,\ldots,G_m$ merupakan himpunan
terukur saling lepas yang mempunyai ukuran berhingga.

(c) Jika $E_0,\ldots,E_n$ merupakan himpunan terukur yang mempunyai ukuran berhingga, dan
$a_0,\ldots,a_n\in\Bbb R$ sedemikian sehingga $\sum_{i=0}^na_i\chi E_i(x)\ge
0$ untuk setiap $x\in X$, maka $\sum_{i=0}^na_i\mu E_i\ge 0$.

\proof{{\bf (a)} Tetapkan $m=2^{n+1}-2$, dan daftarkan
subhimpunan-subhimpunan tak kosong dari $\{0,\ldots,n\}$ sebagai $I_0,\ldots,I_m$. Untuk setiap
$j\le m$, tetapkan

\Centerline{$G_j
=\bigcap_{i\in I_j}E_i\setminus\bigcup_{i\le n,i\notin I_j}E_i$.}

\noindent   Maka setiap $G_j$ merupakan himpunan terukur, karena diperoleh dari
berhingga banyak himpunan terukur melalui operasi $\cup$, $\cap$, dan
$\setminus$, serta mempunyai ukuran berhingga, karena $I_j\ne\emptyset$ dan
$G_j\subseteq E_i$ jika $i\in I_j$. Selain itu, himpunan-himpunan $G_j$ saling lepas,
sebab jika $i\in I_j\setminus I_k$, maka $G_j\subseteq E_i$ dan $G_k\cap
E_i=\emptyset$. Akhirnya, jika $k\le n$ dan $x\in E_k$, terdapat
$j\le
m$ sedemikian sehingga $I_j=\{i:i\le n,\,x\in E_i\}$, dan dalam hal ini $x\in
G_j\subseteq E_k$; dengan demikian, $E_k$ merupakan gabungan dari semua $G_j$ yang
termuat di dalamnya.

\medskip

{\bf (b)} Nyatakan $f$ sebagai $\sum_{i=0}^na_i\chi E_i$, dengan
$E_0,\ldots,E_n$ merupakan himpunan terukur yang mempunyai ukuran berhingga dan
$a_0,\ldots,a_n$ merupakan bilangan real. Misalkan $G_0,\ldots,G_m$ adalah himpunan
terukur saling lepas yang mempunyai ukuran berhingga sedemikian sehingga setiap $E_i$ dapat dinyatakan
sebagai gabungan dari $G_j$ yang sesuai. Tetapkan $c_{ij}=1$ jika $G_j\subseteq
E_i$,
dan $0$ jika tidak, sehingga, karena himpunan-himpunan $G_j$ saling lepas, $\chi
E_i=\sum_{j=0}^mc_{ij}\chi G_j$ untuk setiap $i$. Maka

\Centerline{$f=\sum_{i=0}^na_i\chi E_i
=\sum_{i=0}^n\sum_{j=0}^ma_ic_{ij}\chi G_j
=\sum_{j=0}^mb_j\chi G_j$,}

\noindent dengan menetapkan $b_j=\sum_{i=0}^na_ic_{ij}$ untuk setiap $j\le m$.

\medskip

{\bf (c)} Tetapkan $f=\sum_{i=0}^na_i\chi E_i$, dan ambil $G_j$, $c_{ij}$,
$b_j$ seperti dalam (b). Maka $b_j\mu G_j\ge 0$ untuk setiap $j$. \Prf\ Jika
$G_j=\emptyset$, hal ini jelas. Jika tidak, ambil $x\in G_j$; maka

\Centerline{$0\le f(x)=\sum_{i=0}^mb_i\chi G_i(x)=b_j\chi
G_j(x)=b_j$,}

\noindent maka $b_j\mu G_j\ge 0$ pula.\ \Qed\ Selanjutnya, karena
himpunan-himpunan $G_j$
saling lepas,

\Centerline{$\mu E_i=\sum_{j=0}^mc_{ij}\mu G_j$}

\noindent untuk setiap $i$, sehingga

\Centerline{$\sum_{i=0}^na_i\mu E_i
=\sum_{i=0}^n\sum_{j=0}^ma_ic_{ij}\mu G_j
=\sum_{j=0}^mb_j\mu G_j
\ge 0$,}

\noindent sebagaimana dikehendaki.
}%end of proof of 122C

\leader{122D}{Korolari} Misalkan $(X,\Sigma,\mu)$ suatu ruang ukur. Jika

\Centerline{$\sum_{i=0}^m a_i\chi E_i
=\sum_{j=0}^n b_j\chi F_j$,}

\noindent dengan semua $E_i$ dan $F_j$ merupakan
himpunan terukur yang mempunyai ukuran berhingga dan semua $ a_i$, $ b_j$ merupakan bilangan
real, maka

\Centerline{$\sum_{i=0}^m a_i\mu E_i=\sum_{j=0}^n b_j\mu F_j$.}

\proof{ Terapkan 122Cc pada
$\sum_{i=0}^m a_i\chi E_i+\sum_{j=0}^n(-b_j)\chi F_j$ untuk melihat bahwa
$\sum_{i=0}^ma_i\mu E_i-\sum_{j=0}^nb_j\mu F_j\ge 0$; sekarang pertukarkan
peran
kedua jumlah tersebut untuk memperoleh pertidaksamaan yang berlawanan.
}%end of proof of 122D


\leader{122E}{Definisi} Misalkan $(X,\Sigma,\mu)$ suatu ruang ukur.
Dengan demikian, kita
dapat mendefinisikan {\bf integral} $\int f$ dari $f$, untuk fungsi sederhana
$f:X\to\Bbb R$, melalui ketentuan bahwa
$\int f=\sum_{i=0}^m a_i\mu E_i$ setiap kali $f=\sum_{i=0}^m a_i\chi
E_i$ dan setiap $E_i$ merupakan himpunan terukur yang mempunyai ukuran berhingga\cmmnt{;
122D menjamin bahwa representasi $f$ yang kita pilih
untuk perhitungan tidak akan memengaruhi hasilnya}.

\leader{122F}{Proposisi} Misalkan $(X,\Sigma,\mu)$ suatu ruang ukur.

(a) Jika $f$, $g:X\to\Bbb R$ adalah fungsi sederhana, maka $f+g$ merupakan
fungsi
sederhana dan $\int f+g=\int f+\int g$.

(b) Jika $f$ merupakan fungsi sederhana dan $c\in\Bbb R$, maka $cf$ merupakan
fungsi
sederhana dan $\int cf=c\int f$.

(c) Jika $f$, $g$ merupakan fungsi sederhana dan $f(x)\le g(x)$ untuk setiap
$x\in
X$, maka $\int f\le \int g$.

\proof{ (a) dan (b) langsung mengikuti rumus yang diberikan
untuk $\int f$ dalam 122E. Untuk (c), perhatikan bahwa $g-f$ merupakan fungsi
sederhana nonnegatif, sehingga $\int g-f\ge 0$, menurut 122Cc; tetapi hal ini berarti
bahwa
$\int g-\int f\ge 0$.
}%end of proof of 122F


\vleader{72pt}{122G}{Lema} Misalkan $(X,\Sigma,\mu)$ suatu ruang ukur.
Jika $\langle f_n\rangle_{n\in\Bbb N}$ merupakan barisan fungsi sederhana
yang tak-menurun (dalam arti bahwa $f_n(x)\le f_{n+1}(x)$ untuk
setiap $n\in\Bbb N$, $x\in X$) dan $f$ merupakan fungsi sederhana sedemikian sehingga
$f(x)\le\sup_{n\in\Bbb N}f_n(x)$ untuk hampir setiap $x\in X$ (dengan memperbolehkan
$\sup_{n\in\Bbb N}f_n(x)=\infty$ dalam rumus ini), maka
$\int f\le\sup_{n\in\Bbb N}\int f_n$.

\proof{ Perhatikan bahwa $f-f_0$ merupakan fungsi sederhana, sehingga
$H=\{x:(f-f_0)(x)\ne 0\}$ adalah gabungan berhingga dari himpunan-himpunan yang mempunyai ukuran berhingga,
dan $\mu H<\infty$; selain itu, $f-f_0$ terbatas, sehingga terdapat $M\ge 0$
sedemikian sehingga $(f-f_0)(x)\le M$ untuk setiap $x\in X$.

Misalkan $\epsilon>0$. Untuk setiap $n\in\Bbb N$, tetapkan
$H_n=\{x:(f-f_n)(x)\ge\epsilon\}$. Maka setiap $H_n$ terukur (menurut
121E), dan $\langle H_n\rangle_{n\in\Bbb N}$ merupakan barisan
himpunan yang tak-menaik dengan irisan

\Centerline{$\bigcap_{n\in\Bbb N}H_n
=\{x:f(x)\ge\epsilon+\sup_{n\in\Bbb N}f_n(x)\}
\subseteq\{x:f(x)>\sup_{n\in\Bbb N}f_n(x)\}$.}

\noindent Karena $f(x)\le\sup_{n\in\Bbb N}f_n(x)$ untuk hampir setiap
$x$,
$\{x:f(x)>\sup_{n\in\Bbb N}f_n(x)\}$ dan $\bigcap_{n\in\Bbb N}H_n$
terabaikan. Selain itu, $\mu H_0<\infty$, karena
$H_0\subseteq H$.
Akibatnya,

\Centerline{$\lim_{n\to\infty}\mu H_n=\mu(\bigcap_{n\in\Bbb N}H_n)=0$}

\noindent (112Cf). Ambil $n$ cukup besar sehingga $\mu H_n\le\epsilon$.

Tinjau fungsi sederhana $g=f_n+\epsilon\chi H+M\chi H_n$. Maka
$f\le g$, sehingga

\Centerline{$\int f\le\int g=\int f_n+\epsilon\mu H+M\mu H_n
\le\int f_n+\epsilon(M+\mu H)$.}

\noindent Karena $\epsilon$ sebarang,
$\int f\le\sup_{n\in\Bbb N}\int f_n$.
}%end of proof of 122G

 
\leader{122H}{Definisi} Misalkan $(X,\Sigma,\mu)$ suatu ruang ukur.
Untuk sisa bagian ini, saya akan menuliskan $U$ untuk himpunan
fungsi $f$ sedemikian sehingga

\quad(i) domain $f$ merupakan subhimpunan koterabaikan dari $X$ dan
$f(x)\in\coint{0,\infty}$ untuk setiap $x\in\dom f$,

\quad(ii) terdapat barisan tak-menurun $\sequencen{f_n}$ dari
fungsi sederhana nonnegatif sedemikian sehingga
$\sup_{n\in\Bbb N}\int f_n<\infty$ dan
$\lim_{n\to\infty}f_n(x)\penalty-100=f(x)$ untuk hampir
setiap $x\in X$.

\leader{122I}{Lema} Jika $f$ dan $\langle f_n\rangle_{n\in\Bbb N}$
seperti dalam 122H, maka

\Centerline{$\sup_{n\in\Bbb N}\int f_n
=\sup\{\int g:g$ adalah fungsi sederhana dan $g\leae f\}$.}

\proof{ Tentu saja

\Centerline{$\sup_{n\in\Bbb N}\int f_n
\le \sup\{\int g:g$ adalah fungsi sederhana dan $g\leae f\}$}

\noindent karena $f_n\leae f$ untuk setiap $n$. Di pihak lain, jika
$g$ adalah fungsi sederhana dan $g\leae f$, maka
$g(x)\le\sup_{n\in\Bbb N}f_n(x)$ untuk hampir setiap $x$, sehingga
$\int g\le\sup_{n\in\Bbb N}\int f_n$ menurut 122G. Dengan demikian

\Centerline{$\sup_{n\in\Bbb N}\int f_n
\ge \sup\{\int g:g$ adalah fungsi sederhana dan $g\leae f\}$,}

\noindent sebagaimana diperlukan.
}%end of proof of 122I


\leader{122J}{Lema} Misalkan $(X,\Sigma,\mu)$ suatu ruang
ukur\cmmnt{, dan definisikan $U$ seperti dalam 122H}.

(a) Jika $f$ adalah fungsi yang didefinisikan pada subhimpunan
koterabaikan dari $X$ dan mengambil nilai dalam $\coint{0,\infty}$,
maka $f\in U$ jika dan hanya jika terdapat himpunan terukur
koterabaikan $E\subseteq\dom f$ sedemikian sehingga

\inset{($\alpha$)
$f\restr E$ terukur,

($\beta$) untuk setiap $\epsilon>0$,
$\mu\{x:x\in E,\,f(x)\ge\epsilon\}<\infty$,

($\gamma$)
$\sup\{\int g:g$ adalah fungsi sederhana, $g\leae f\}<\infty$.}

(b) Andaikan bahwa $f\in U$ dan bahwa $h$ adalah fungsi yang
didefinisikan pada subhimpunan koterabaikan dari $X$ dan mengambil
nilai dalam $\coint{0,\infty}$. Andaikan bahwa $h\leae f$ dan terdapat
himpunan koterabaikan $F\subseteq X$ sedemikian sehingga $h\restr F$
terukur. Maka $h\in U$.

\proof{{\bf (a)(i)} Andaikan bahwa $f\in U$. Maka terdapat barisan
tak-menurun $\langle f_n\rangle_{n\in\Bbb N}$ dari fungsi sederhana
nonnegatif sedemikian sehingga $f\eae\lim_{n\to\infty}f_n$ dan
$\sup_{n\in\Bbb N}\int f_n=c<\infty$. Himpunan
$\{x:f(x)=\lim_{n\to\infty}f_n(x)\}$ koterabaikan, sehingga memuat
suatu himpunan terukur koterabaikan, sebut saja $E$. Sekarang
$f\restr E=(\lim_{n\to\infty}f_n)\restr E$ terukur, menurut 121Fa
dan 121Eh; dengan demikian ($\alpha$) terpenuhi. Selanjutnya, untuk
$\epsilon>0$, tetapkan
$H_n=\{x:x\in E,\,f_n(x)\ge{1\over 2}\epsilon\}$; maka
$f_n\ge{1\over 2}\epsilon\chi H_n$, sehingga

\Centerline{$\Bover12\epsilon\mu H_n
=\int\bover12\epsilon\chi H_n\le\int f_n\le c$,}

\noindent untuk setiap $n$. Sekarang $\langle H_n\rangle_{n\in\Bbb N}$
merupakan barisan tak-menurun, sehingga

\Centerline{$\mu(\bigcup_{n\in\Bbb N}H_n)
=\sup_{n\in\Bbb N}\mu H_n\le 2c/\epsilon$,}

\noindent menurut 112Ce. Oleh karena itu

\Centerline{$\mu\{x:x\in E,\,f(x)\ge\epsilon\}
\le\mu(\bigcup_{n\in\Bbb N}H_n)\le 2c/\epsilon<\infty$.}

\noindent Karena $\epsilon$ sebarang, ($\beta$) terpenuhi.
Terakhir, ($\gamma$) terpenuhi menurut 122I.

\medskip

\quad{\bf (ii)} Sekarang andaikan bahwa syarat-syarat
($\alpha$)-($\gamma$) terpenuhi. Ambil himpunan koterabaikan
$E\in\Sigma$ yang sesuai, dan untuk setiap $n\in\Bbb N$ definisikan
$f_n:X\to\Bbb R$ dengan menetapkan

$$\eqalign{f_n(x)&=2^{-n}k\text{ jika }x\in E,\, 0\le k<4^n,\,
2^{-n}k\le f(x)<2^{-n}(k+1),\cr
&=0\text{ jika }x\in X\setminus E,\cr
&=2^n\text{ jika }x\in E\text{ dan }f(x)\ge 2^n.\cr}$$

\noindent Maka $f_n$ merupakan fungsi sederhana nonnegatif, karena
dapat dinyatakan sebagai

\Centerline{$f_n
=\sum_{k=1}^{4^n}2^{-n}\chi\{x:x\in E,\,f(x)\ge 2^{-n}k\}$;}

\noindent semua himpunan $\{x:x\in E,\,f(x)\ge 2^{-n}k\}$ terukur
(karena $f\restr E$ terukur) dan mempunyai ukuran berhingga, menurut
($\beta$).
Juga mudah dilihat bahwa $\langle f_n\rangle_{n\in\Bbb N}$ merupakan
barisan tak-menurun yang konvergen ke $f$ pada setiap titik di $E$,
sehingga $f\eae\lim_{n\to\infty}f_n$. Terakhir,

\Centerline{$\lim_{n\to\infty}\int f_n
=\sup_{n\in\Bbb N}\int f_n
\le\sup\{\int g:g\text{ adalah fungsi sederhana},\,g\leae f\}
<\infty$,}

\noindent menurut ($\gamma$). Jadi $f\in U$.

\medskip

{\bf (b)} Misalkan $E$ suatu himpunan seperti dalam (a). Himpunan
$E$, $F$, dan $\{x:h(x)\le f(x)\}$ semuanya koterabaikan, sehingga
terdapat himpunan terukur koterabaikan $E'$ yang termuat dalam irisan
ketiganya. Sekarang $E'\subseteq\dom h$, $h\restr E'$ terukur,

\Centerline{$\mu\{x:x\in E',\,h(x)\ge\epsilon\}
\le\mu\{x:x\in E,\,f(x)\ge\epsilon\}<\infty$}

\noindent untuk setiap $\epsilon>0$, dan

\Centerline{$\sup\{\int g:g\text{ adalah fungsi sederhana},\, g\leae h\}
\le\sup\{\int g:g\text{ adalah fungsi sederhana},\,g\leae f\}<\infty$.}

\noindent Jadi $h\in U$.
}%end of proof of 122J

\leader{122K}{Definisi} Misalkan $(X,\Sigma,\mu)$ suatu ruang
ukur\cmmnt{, dan definisikan $U$ seperti dalam 122H}. Untuk $f\in U$,
tetapkan

\Centerline{$\int f
=\sup\{\int g:g$ adalah fungsi sederhana dan $g\leae f\}$.}

\cmmnt{\noindent Menurut 122I, kita melihat bahwa
$\int f=\lim_{n\to\infty}\int f_n$ setiap kali $\sequencen{f_n}$
merupakan barisan tak-menurun dari fungsi sederhana yang konvergen ke
$f$ hampir di mana-mana; khususnya, jika $f\in U$ sendiri merupakan
fungsi sederhana, maka $\int f$, sebagaimana didefinisikan di sini,
sesuai dengan definisi semula untuk $\int f$ dalam 122E, karena kita
dapat mengambil $f_n=f$ untuk setiap $n$.}



\leader{122L}{Lema} Misalkan $(X,\Sigma,\mu)$ suatu ruang ukur.

(a) Jika $f$, $g\in U$, maka $f+g\in U$ dan
$\int f+g=\int f+\int g$.

(b) Jika $f\in U$ dan $c\ge 0$, maka $cf\in U$ dan
$\int cf=c\int f$.

(c) Jika $f$, $g\in U$ dan $f\leae g$, maka $\int f\le \int g$.

(d) Jika $f\in U$ dan $g$ adalah fungsi dengan domain berupa
subhimpunan koterabaikan dari $X$, mengambil nilai dalam
$\coint{0,\infty}$, dan sama dengan $f$ hampir di mana-mana, maka
$g\in U$ dan $\int g=\int f$.

(e) Jika $f_1$, $g_1$, $f_2$, $g_2\in U$ dan
$f_1-f_2=g_1-g_2$, maka
$\int f_1-\int f_2=\int g_1-\int g_2$.

\proof{{\bf (a)} Kita tahu bahwa terdapat barisan-barisan tak-menurun
$\langle f_n\rangle_{n\in\Bbb N}$,
$\langle g_n\rangle_{n\in\Bbb N}$ dari fungsi sederhana nonnegatif
sedemikian sehingga $f\eae\lim_{n\to\infty}f_n$,
$g\eae\lim_{n\to\infty}g_n$, $\sup_{n\in\Bbb N}\int f_n<\infty$,
dan $\sup_{n\in\Bbb N}\int g_n<\infty$. Sekarang
$\langle f_n+g_n\rangle_{n\in\Bbb N}$ merupakan barisan tak-menurun
dari fungsi sederhana yang konvergen ke $f+g$ hampir di mana-mana, dan

\Centerline{$\sup_{n\in\Bbb N}\int f_n+g_n
=\lim_{n\to\infty}\int f_n+g_n
=\lim_{n\to\infty}\int f_n+\lim_{n\to\infty}\int g_n=\int f+\int g$.}

\noindent Oleh karena itu $f+g\in U$, dan juga, sebagaimana dicatat
dalam 122K,

\Centerline{$\int f+g=\lim_{n\to\infty}\int f_n+g_n=\int f+\int g$.}

\medskip

{\bf (b)} Kita tahu bahwa terdapat barisan tak-menurun
$\langle f_n\rangle_{n\in\Bbb N}$ dari fungsi sederhana nonnegatif
sedemikian sehingga $f\eae\lim_{n\to\infty}f_n$ dan
$\sup_{n\in\Bbb N}\int f_n<\infty$. Sekarang
$\langle cf_n\rangle_{n\in\Bbb N}$ merupakan barisan tak-menurun dari
fungsi sederhana yang konvergen ke $cf$ hampir di mana-mana, dan

\Centerline{$\sup_{n\in\Bbb N}\int cf_n
=\lim_{n\to\infty}\int cf_n=c\lim_{n\to\infty}\int f_n=c\int f$.}

\noindent Oleh karena itu $cf\in U$, dan juga, sebagaimana dicatat
dalam 122K,

\Centerline{$\int cf=\lim_{n\to\infty}\int cf_n=c\int f$.}

\medskip

{\bf (c)} Hal ini langsung dari 122K.

\medskip

{\bf (d)} Jika $\langle f_n\rangle_{n\in\Bbb N}$ merupakan barisan
tak-menurun dari fungsi sederhana yang konvergen ke $f$ hampir di
mana-mana, maka barisan itu juga konvergen ke $g$ hampir di mana-mana;
jadi $g\in U$ dan

\Centerline{$\int g=\lim_{n\to\infty}\int f_n=\int f$.}

\medskip

{\bf (e)} Menurut (a), $f_1+g_2$ dan $f_2+g_1$ keduanya termasuk
dalam $U$. Selain itu, keduanya sama pada setiap titik tempat keempat
fungsi terdefinisi, dan ini berlaku hampir di mana-mana. Jadi

\Centerline{$\int f_1+\int g_2=\int f_1+g_2
=\int f_2+g_1=\int f_2+\int g_1$,}

\noindent dengan menggunakan (a) dan (d). Dengan memindahkan
$\int g_2$ dan $\int f_2$ ke ruas lain persamaan, kita memperoleh
hasilnya.
}%end of proof of 122L
\leader{122M}{Definisi} Misalkan $(X,\Sigma,\mu)$ suatu ruang ukur.
\cmmnt{Definisikan $U$ seperti dalam 122H.} Fungsi bernilai real $f$ disebut {\bf
terintegralkan}, atau {\bf terintegralkan pada
$X$}, atau {\bf $\mu$-terintegralkan pada $X$}, jika dapat dinyatakan sebagai
$f_1-f_2$ dengan $f_1$, $f_2\in U$, dan dalam hal ini {\bf integralnya}
adalah

\Centerline{$\int f=\int f_1-\int f_2$.}

\cmmnt{
\leader{122N}{Catatan (a)} Dari 122Le kita melihat bahwa
integral
$\int f$ didefinisikan secara tunggal oleh
rumus dalam 122M. Kedua, jika $f\in U$, maka $f=f-\tbf{0}$
terintegralkan, dan
integralnya di sini sama dengan integral yang didefinisikan dalam 122K. Terakhir, jika $f$
adalah fungsi sederhana, maka fungsi tersebut dapat dinyatakan sebagai $f_1-f_2$ dengan $f_1$,
$f_2$ fungsi sederhana nonnegatif (jika $f=\sum_{i=0}^na_i\chi
E_i$, dengan setiap $E_i$ terukur dan berukuran berhingga, tetapkan

\Centerline{$f_1=\sum_{i=0}^na_i^+\chi E_i$,
\quad$f_2=\sum_{i=0}^na_i^-\chi E_i$,}

\noindent dengan menuliskan $a_i^+=\max(a_i,0)$, $a_i^-=\max(-a_i,0)$); sehingga

\Centerline{$\int f=\int f_1-\int f_2=$$\sum_{i=0}^na_i\mu E_i$,}

\noindent dan definisi dalam 122M selaras dengan definisi
dalam 122E.


\header{122Nb}{\bf (b)} Notasi-notasi alternatif yang akan saya gunakan untuk
$\int
f$ ialah $\int_Xf$,
$\int fd\mu$, $\int f(x)\mu(dx)$, $\int f(x)dx$, $\int_Xf(x)\mu(dx)$,
dan sebagainya, bergantung pada aspek konteks mana yang tampaknya perlu ditekankan.

Jika $\mu$ adalah ukuran Lebesgue pada $\Bbb R$ atau $\BbbR^r$, kita mengatakan bahwa
$\int f$ adalah {\bf integral Lebesgue} dari $f$, dan bahwa $f$ {\bf
terintegralkan secara Lebesgue} jika integral ini terdefinisi.

\header{122Nc}{\bf (c)} Perhatikan bahwa ketika saya mengatakan, dalam 122M, bahwa `$f$ dapat
dinyatakan
sebagai $f_1-f_2$', yang saya maksud ialah
menafsirkan $f_1-f_2$ menurut aturan-aturan yang ditetapkan dalam 121E, sehingga
$\dom f$ harus sama dengan $\dom(f_1-f_2)=\dom f_1\cap\dom f_2$, dan tentu saja
koterabaikan.
}%end of comment

\leader{122O}{Teorema} Misalkan $(X,\Sigma,\mu)$ suatu ruang ukur.

(a) Jika $f$ dan $g$ terintegralkan pada $X$,
maka $f+g$ terintegralkan dan $\int f+g=\int f+\int g$.

(b) Jika $f$ terintegralkan pada $X$ dan $c\in\Bbb R$, maka $cf$
terintegralkan dan $\int cf=c\int f$.

(c) Jika $f$ terintegralkan pada $X$ dan $f\ge 0$ a.e.\ maka
$\int f\ge 0$.

(d) Jika $f$ dan $g$ terintegralkan pada $X$ dan $f\leae g$,
maka $\int f\le\int g$.

\proof{{\bf (a)} Nyatakan $f$ sebagai $f_1-f_2$ dan $g$ sebagai $g_1-g_2$
dengan $f_1$, $f_2$, $g_1$, dan $g_2$ termasuk dalam $U$, sebagaimana didefinisikan dalam 122H.
Maka $f+g=(f_1+g_1)-(f_2+g_2)$ terintegralkan karena $U$ tertutup
terhadap penjumlahan (122La), dan

\Centerline{$\int f+g=\int f_1+g_1-\int f_2+g_2
=\int f_1+\int g_1-\int f_2-\int g_2=\int f+\int g$.}

\medskip

{\bf (b)} Nyatakan $f$ sebagai $f_1-f_2$ dengan $f_1$, $f_2$ termasuk dalam $U$.
Jika $c\ge 0$, maka $cf=cf_1-cf_2$ terintegralkan karena $U$ tertutup
terhadap perkalian dengan skalar nonnegatif (122Lb), dan

\Centerline{$\int cf=\int cf_1-\int cf_2=c\int f_1-c\int f_2=c\int
f$.}

\noindent Jika $c=-1$, maka $-f=f_2-f_1$ terintegralkan dan

\Centerline{$\int (-f)=\int f_2-\int f_1=-\int f$.}

\noindent Dengan menggabungkan
kedua hal ini, kita memperoleh hasil untuk $c<0$.

\medskip

{\bf (c)} Nyatakan $f$ sebagai $f_1-f_2$ dengan $f_1$, $f_2\in U$. Maka
$f_2\leae f_1$, sehingga $\int f_2\le\int f_1$ (122Lc), dan $\int f\ge 0$.

\medskip

{\bf (d)} Terapkan (c) pada $g-f$.
}%end of proof of 122O


\leader{122P}{Teorema} Misalkan $(X,\Sigma,\mu)$ suatu ruang ukur dan
$f$ suatu
fungsi bernilai real yang didefinisikan pada subhimpunan koterabaikan dari $X$. Maka
syarat-syarat berikut ekuivalen:

\quad(i) $f$ terintegralkan;

\quad(ii) $|f|\in U$\cmmnt{, sebagaimana didefinisikan dalam 122H,} dan terdapat
himpunan koterabaikan $E\subseteq X$ sedemikian sehingga $f\restr E$ terukur;

\quad(iii) terdapat $g\in U$ dan
himpunan koterabaikan $E\subseteq X$ sedemikian sehingga $|f|\leae g$
dan $f\restr E$ terukur.

\proof{{\bf (i)$\Rightarrow$(iii)} Andaikan bahwa $f$
terintegralkan. Misalkan $f_1$,
$f_2\in U$ sedemikian sehingga $f=f_1-f_2$. Maka terdapat himpunan-himpunan koterabaikan
$E_1$, $E_2$ sedemikian sehingga $f_1\restr E_1$ dan $f_2\restr
E_2$ terukur; tetapkan $E=E_1\cap E_2$, sehingga $E$ juga merupakan
himpunan koterabaikan. Sekarang $f\restr E=f_1\restr E-f_2\restr E$
terukur. Selanjutnya, $f_1+f_2\in U$ (122La) dan $|f|(x)\le
f_1(x)+f_2(x)$ untuk setiap $x\in\dom f$, sehingga kita dapat mengambil $g=f_1+f_2$.

\medskip

{\bf (iii)$\Rightarrow$(ii)} Jika $f\restr E$ terukur, maka
$|f|\restr E=|f\restr E|$ juga terukur (121Eg); jadi, jika $g\in U$ dan
$|f|\leae g$, maka $|f|\in U$ menurut 122Jb.

\medskip

{\bf (ii)$\Rightarrow$(i)} Andaikan bahwa $f$ memenuhi syarat-syarat
dalam
(ii). Tetapkan
$f^+={1\over 2}(|f|+f)$ dan $f^-={1\over 2}(|f|-f)$. Tentu saja,
$|f|\restr E$, $f^+\restr E$, dan $f^-\restr E$ semuanya terukur.
Selain itu, $0\le f^+(x)\le |f|(x)$ dan $0\le f^-(x)\le |f|(x)$ untuk setiap
$x\in\dom f$, sedangkan $|f|\in U$ menurut hipotesis, sehingga $f^+$ dan $f^-$
termasuk
dalam $U$ menurut 122Jb. Akhirnya, $f=f^+-f^-$, sehingga $f$ terintegralkan.
}%end of proof of 122P

\leader{122Q}{Catatan} Syarat `terdapat himpunan koterabaikan $E$
sedemikian sehingga $f\restr E$ terukur' muncul begitu sering sehingga menurut saya
patut diberi sebuah istilah; fungsi-fungsi demikian akan saya sebut {\bf
terukur secara
virtual}, atau {\bf $\mu$-terukur secara virtual} jika ukuran yang dimaksud
tampaknya perlu dinyatakan.

\leader{122R}{Korolari} Misalkan $(X,\Sigma,\mu)$ suatu ruang ukur.

(a) Suatu fungsi bernilai real nonnegatif, yang didefinisikan pada subhimpunan $X$,
terintegralkan jika dan hanya jika fungsi itu termasuk dalam $U$\cmmnt{, sebagaimana didefinisikan dalam 122H}.

(b) Jika $f$ terintegralkan pada $X$ dan $h$ suatu fungsi bernilai real
yang didefinisikan pada subhimpunan koterabaikan dari $X$ dan sama dengan $f$ hampir
di mana-mana, maka $h$ terintegralkan, dengan $\int h=\int f$.

(c) Jika $f$ terintegralkan pada $X$, $f\ge 0$ a.e.\ dan $\int f\le 0$,
maka $f=0$ a.e.

(d) Jika $f$ dan $g$ terintegralkan pada $X$, $f\leae g$, dan
$\int g\le\int f$, maka $f\eae g$.

(e) Jika $f$ terintegralkan pada $X$, maka $|f|$ juga terintegralkan, dan
$|\int f|\le\int|f|$.

\proof{{\bf (a)} Jika $f$ terintegralkan, maka $f=|f|\in U$, menurut
122P(ii). Jika $f\in U$, maka $f=f-\tbf{0}$ terintegralkan, menurut 122M.

\medskip

{\bf (b)} Misalkan $E$, $F$ himpunan-himpunan koterabaikan sedemikian sehingga $f\restr E$
terukur dan $h\restr F=f\restr F$; maka $E\cap F$ koterabaikan
dan $h\restr(E\cap F)=(f\restr E)\restr F$ terukur. Selanjutnya,
terdapat
$g\in U$ sedemikian sehingga $|f|\leae g$, dan tentu saja $|h|\leae g$.
Jadi, $h$ terintegralkan menurut 122P(iii). Dengan menerapkan 122Od pada $f$ dan $h$,
lalu pada $h$ dan $f$, diperoleh $\int h=\int f$.

\medskip

{\bf (c)} \Quer\ Andaikan, jika mungkin, bahwa yang sebaliknya berlaku. Misalkan $E\subseteq X$
himpunan koterabaikan sedemikian sehingga $f\restr E$ terukur (122P(ii)),
dan ambil $E'\subseteq E\cap\dom f$ suatu himpunan terukur koterabaikan. Maka
$F=\{x:x\in E',\,f(x)>0\}$ pasti tidak terabaikan. Tetapkan
$F_k=\{x:x\in
E',\,f(x)\ge 2^{-k}\}$ untuk setiap $k\in\Bbb N$, sehingga
$F=\bigcup_{k\in\Bbb N}F_k$ dan terdapat $k$ sedemikian sehingga $\mu F_k>0$.
Namun, tinjau $g=2^{-k}\chi F_k$. Karena $f\ge 0$ a.e.\ dan
$f\ge 2^{-k}$ pada $F_k$, berlaku $f\geae g$, sehingga

\Centerline{$ 0<2^{-k}\mu F_k=\int g\le\int f$,}

\noindent menurut 122Od; hal ini mustahil. \Bang

\medskip

{\bf (d)} Terapkan (c) pada $g-f$.

\medskip

{\bf (e)} Menurut (i)$\Rightarrow$(ii) dalam 122P, $|f|$ terintegralkan. Sekarang
$f^+={1\over 2}(|f|+f)$ dan $f^-={1\over 2}(|f|-f)$ keduanya
terintegralkan
dan nonnegatif, sehingga integralnya nonnegatif, dan

\Centerline{$|\int f|=|\int f^+-\int f^-|\le\int f^++\int
f^-=\int|f|$.}
}%end of proof of 122R
 
\exercises{
\leader{122X}{Latihan dasar (a)} Misalkan $(X,\Sigma,\mu)$ suatu ruang
ukur. (i) Tunjukkan bahwa jika $f:X\to\Bbb R$ adalah fungsi sederhana,
maka $|f|$ juga sederhana, dengan menetapkan $|f|(x)=|f(x)|$ untuk
$x\in\dom f=X$. (ii) Tunjukkan bahwa jika $f$, $g:X\to\Bbb R$ adalah
fungsi-fungsi sederhana, maka $f\vee g$ dan $f\wedge g$, sebagaimana
didefinisikan dalam 121Xb, juga sederhana.
%122D

\sqheader 122Xb Misalkan $(X,\Sigma,\mu)$ suatu ruang ukur dan $f$
suatu fungsi bernilai real yang terintegralkan pada $X$. Tunjukkan bahwa
untuk setiap $\epsilon>0$ terdapat fungsi sederhana $g:X\to\Bbb R$
sedemikian sehingga $\int|f-g|\le\epsilon$. \Hint{tinjau terlebih dahulu
$f$ yang nonnegatif.}
%122O

\spheader 122Xc Misalkan $(X,\Sigma,\mu)$ suatu ruang ukur, dan tuliskan
$\eusm L^1$ untuk himpunan semua fungsi bernilai real yang terintegralkan
atas $X$. Misalkan $\Phi\subseteq\eusm L^1$ sedemikian sehingga

\inset{(i) $\chi E\in\Phi$ setiap kali $E\in\Sigma$ dan
$\mu E<\infty$;

(ii) $f+g\in\Phi$ untuk semua $f$, $g\in\Phi$;

(iii) $cf\in\Phi$ setiap kali $c\in\Bbb R$, $f\in\Phi$;

(iv) $f\in\Phi$ setiap kali $f\in\eusm L^1$ sedemikian sehingga terdapat
barisan tak-menurun $\sequencen{f_n}$ di dalam $\Phi$ dengan
$\lim_{n\to\infty}f_n=f$ hampir di mana-mana.}

\noindent Tunjukkan bahwa $\Phi=\eusm L^1$.
%122O

\sqheader 122Xd Misalkan $\mu$ ukuran cacah pada $\Bbb N$ (112Bd).
Tunjukkan bahwa suatu fungsi $f:\Bbb N\to\Bbb R$ (yakni, suatu barisan
$\sequencen{f(n)}$) terintegralkan terhadap $\mu$ jika dan hanya jika
barisan itu dapat dijumlahkan secara absolut, dan dalam hal ini

\Centerline{$\int fd\mu=\int_{\Bbb N}f(n)\mu(dn)
=\sum_{n=0}^{\infty}f(n)$.}
%122O

\sqheader 122Xe Misalkan $(X,\Sigma,\mu)$ suatu ruang ukur dan $f$, $g$
dua fungsi bernilai real yang terukur secara virtual dan didefinisikan
pada subhimpunan-subhimpunan $X$. (i) Tunjukkan bahwa $f+g$, $f\times g$
dan $f/g$, yang didefinisikan sebagaimana dalam 121E, semuanya terukur
secara virtual. (ii) Tunjukkan bahwa jika $h$ suatu fungsi bernilai real
terukur Borel yang didefinisikan pada sembarang subhimpunan $\Bbb R$,
maka komposisi $hf$ terukur secara virtual.
%122Q

\sqheader 122Xf Misalkan $(X,\Sigma,\mu)$ suatu ruang ukur dan
$\sequencen{f_n}$ suatu barisan fungsi bernilai real yang terukur secara
virtual dan didefinisikan pada subhimpunan-subhimpunan $X$. Tunjukkan
bahwa $\lim_{n\to\infty}f_n$, $\sup_{n\in\Bbb N}f_n$,
$\inf_{n\in\Bbb N}f_n$, $\limsup_{n\to\infty}f_n$ dan
$\liminf_{n\to\infty}f_n$, yang didefinisikan sebagaimana dalam 121F,
terukur secara virtual.
%122Q

\sqheader 122Xg Misalkan $(X,\Sigma,\mu)$ suatu ruang ukur dan $f$, $g$
fungsi-fungsi bernilai real yang terintegralkan pada $X$. Tunjukkan bahwa
$f\wedge g$ dan $f\vee g$, yang didefinisikan sebagaimana dalam 121Xb,
terintegralkan.
%122R

\sqheader 122Xh Misalkan $(X,\Sigma,\mu)$ suatu ruang ukur, $f$ suatu
fungsi bernilai real yang terintegralkan pada $X$, dan $g$ suatu fungsi
bernilai real terbatas yang terukur secara virtual serta didefinisikan
pada subhimpunan koterabaikan dari $X$. Tunjukkan bahwa $f\times g$,
yang didefinisikan sebagaimana dalam 121Ed, terintegralkan.
%122R

%maybe better in 122Y?
\spheader 122Xi Misalkan $X$ suatu himpunan, $\Sigma$ suatu
aljabar-$\sigma$ atas subhimpunan-subhimpunan $X$, dan $\mu_1$, $\mu_2$
dua ukuran dengan domain $\Sigma$. Tetapkan
$\mu E=\mu_1E+\mu_2E$ untuk $E\in\Sigma$. Tunjukkan bahwa untuk setiap
fungsi bernilai real $f$ yang didefinisikan pada suatu subhimpunan $X$,
$\int fd\mu=\int fd\mu_1+\int fd\mu_2$ dalam arti bahwa jika salah satu
ruas terdefinisi sebagai bilangan real, maka ruas lainnya juga demikian,
dan keduanya kemudian sama. \Hint{($\alpha$) Periksa bahwa suatu
subhimpunan $X$ bersifat $\mu$-koterabaikan jika dan hanya jika ia
bersifat $\mu_i$-koterabaikan untuk kedua nilai $i$. ($\beta$) Periksa hasil
tersebut untuk fungsi sederhana $f$. ($\gamma$) Sekarang tinjau $f$
nonnegatif yang umum.}

\leader{122Y}{Latihan lanjutan (a)} Misalkan $(X,\Sigma,\mu)$ suatu
ruang ukur `lengkap', yakni ruang yang semua himpunan terabaikannya
terukur (lihat, misalnya, 113Xa). Tunjukkan bahwa jika $f$ suatu fungsi
bernilai real yang terukur secara virtual dan didefinisikan pada suatu
subhimpunan $X$, maka $f$ terukur. Gunakan fakta ini untuk menemukan
penyederhanaan yang sesuai bagi 122J dan 122P untuk ruang-ruang
$(X,\Sigma,\mu)$ semacam itu.
%122P

\spheader 122Yb Tuliskan $\eusm L^1$ untuk himpunan semua fungsi
bernilai real yang terintegralkan secara Lebesgue pada $\Bbb R$. Misalkan
$\Phi\subseteq\eusm L^1$ sedemikian sehingga

\inset{(i) $\chi I\in\Phi$ setiap kali $I$ suatu interval setengah
terbuka terbatas di $\Bbb R$;

(ii) $f+g\in\Phi$ untuk semua $f$, $g\in\Phi$;

(iii) $cf\in\Phi$ setiap kali $c\in\Bbb R$, $f\in\Phi$;

(iv) $f\in\Phi$ setiap kali $f\in\eusm L^1$ sedemikian sehingga terdapat
barisan tak-menurun $\sequencen{f_n}$ di dalam $\Phi$ dengan
$\lim_{n\to\infty}f_n=f$ hampir di mana-mana.}

\noindent Tunjukkan bahwa $\Phi=\eusm L^1$. \Hint{tunjukkan bahwa
($\alpha$) $\chi E\in\Phi$ setiap kali $E$ dapat dinyatakan sebagai
gabungan sejumlah berhingga interval setengah terbuka ($\beta$)
$\chi E\in\Phi$ setiap kali $E$ mempunyai ukuran hingga dan dapat dinyatakan
sebagai gabungan suatu barisan interval setengah terbuka ($\gamma$)
$\chi E\in\Phi$ setiap kali $E$ terukur dan mempunyai ukuran hingga.}
%122O, 122Xc

\spheader 122Yc Misalkan $X$ sembarang himpunan, dan misalkan $\mu$
ukuran cacah pada $X$. Misalkan $f:X\to\Bbb R$ suatu fungsi; tetapkan
$f^+(x)=\max(0,f(x))$, $f^-(x)=\max(0,-f(x))$ untuk $x\in X$.
Tunjukkan bahwa pernyataan-pernyataan berikut ekuivalen: (i)
$\int fd\mu$ ada di $\Bbb R$, dan sama dengan $s$; (ii) untuk setiap
$\epsilon>0$ terdapat $K\subseteq X$ yang berhingga sedemikian sehingga
$|s-\sum_{i\in I}f(i)|\le\epsilon$ setiap kali $I\subseteq X$ adalah
himpunan berhingga yang memuat $K$; (iii) $\sum_{x\in X}f^+(x)$ dan
$\sum_{x\in X}f^-(x)$, yang didefinisikan sebagaimana dalam 112Bd,
berhingga, dan
$s=\sum_{x\in X}f^+(x)-\sum_{x\in X}f^-(x)$.
%122Xd, 122O

\spheader 122Yd Misalkan $(X,\Sigma,\mu)$ suatu ruang ukur. Kita katakan
bahwa suatu fungsi $g:X\to\Bbb R$ bersifat {\bf kuasi-sederhana} jika
dapat dinyatakan sebagai $\sum_{i=0}^{\infty}a_i\chi G_i$, dengan
$\sequence{i}{G_i}$ suatu partisi $X$ menjadi himpunan-himpunan terukur,
$\sequence{i}{a_i}$ suatu barisan di $\Bbb R$, dan
$\sum_{i=0}^{\infty}|a_i|\mu G_i<\infty$, dengan menetapkan
$0\cdot\infty$ sebagai $0$, sehingga boleh terdapat $G_i$ berukuran tak
hingga asalkan $a_i$ yang bersesuaian bernilai nol.

\quad(i) Tunjukkan bahwa jika $g$ dan $h$ adalah fungsi-fungsi
kuasi-sederhana, maka $g+h$, $|g|$ dan $cg$, untuk sembarang
$c\in \Bbb R$, juga kuasi-sederhana. \Hint{untuk $g+h$ Anda akan
memerlukan 111F(b-ii) atau suatu pernyataan yang ekuivalen.}

\quad(ii) Tunjukkan dari prinsip-prinsip pertama (maksud saya, tanpa
menggunakan apa pun yang muncul sesudah 122F dalam bab ini) bahwa jika
$g=\sum_{i=0}^{\infty}a_i\chi G_i$ dan
$h=\sum_{i=0}^{\infty}b_i\chi H_i$ adalah fungsi-fungsi kuasi-sederhana,
dan $g\leae h$, maka
$\sum_{i=0}^{\infty}a_i\mu G_i\le\sum_{i=0}^{\infty}b_i\mu H_i$.

\quad{(iii)} Selanjutnya tunjukkan bahwa kita mempunyai suatu
fungsional $I$ yang didefinisikan dengan menetapkan
$I(g)=\sum_{i=0}^{\infty}a_i\mu G_i$ setiap kali $g$ suatu fungsi
kuasi-sederhana yang direpresentasikan sebagai
$\sum_{i=0}^{\infty}a_i\chi G_i$.

\quad{(iv)} Tunjukkan bahwa jika $g$ dan $h$ adalah fungsi-fungsi
kuasi-sederhana dan $c\in\Bbb R$, maka $I(g+h)=I(g)+I(h)$ dan
$I(cg)=cI(g)$, serta $I(g)\le I(h)$ jika $g\leae h$.

\quad{(v)} Tunjukkan bahwa jika $g$ suatu fungsi kuasi-sederhana, maka
$g$ terintegralkan dan $\int g=I(g)$. (Sekarang saya mengizinkan Anda
menggunakan 122G-122R.)

\quad{(vi)} Tunjukkan bahwa suatu fungsi bernilai real $f$, yang
didefinisikan pada subhimpunan koterabaikan dari $X$, terintegralkan jika
dan hanya jika untuk setiap $\epsilon>0$ terdapat fungsi-fungsi
kuasi-sederhana $g$, $h$ sedemikian sehingga $g\leae f\leae h$ dan
$I(h)-I(g)\le\epsilon$.
%122R

\spheader 122Ye Misalkan $\mu$ ukuran Lebesgue pada $\Bbb R$. Kita
katakan (hanya untuk latihan ini) bahwa suatu fungsi bernilai real $g$
dengan $\dom g\subseteq\Bbb R$ bersifat `pseudo-sederhana' jika dapat
dinyatakan sebagai $\sum_{i=0}^{\infty}a_i\chi J_i$, dengan
$\sequence{i}{J_i}$ suatu barisan interval setengah terbuka terbatas
({\it tidak} harus saling lepas) dan
$\sum_{i=0}^{\infty}|a_i|\mu J_i<\infty$. (Tafsirkan jumlah tak hingga
$\sum_{i=0}^{\infty}a_i\chi J_i$ sebagaimana dalam 121F, sehingga

\Centerline{$\dom(\sum_{i=0}^{\infty}a_i\chi J_i)
=\{x:\lim_{n\to\infty}\sum_{i=0}^na_i(\chi J_i)(x)
  \text{ ada di }\Bbb R\}$.)}

\quad(i) Tunjukkan bahwa jika $g$, $h$ adalah fungsi-fungsi
pseudo-sederhana, maka $g+h$ dan $cg$, untuk sembarang $c\in \Bbb R$,
juga pseudo-sederhana.

\quad(ii) Tunjukkan bahwa jika $g$ suatu fungsi pseudo-sederhana, maka
$g$ terintegralkan.

\quad(iii) Tunjukkan bahwa suatu fungsi bernilai real $f$, yang
didefinisikan pada subhimpunan koterabaikan dari $\Bbb R$,
terintegralkan jika dan hanya jika untuk setiap $\epsilon>0$ terdapat
fungsi-fungsi pseudo-sederhana $g$, $h$ sedemikian sehingga
$g\leae f\leae h$ dan $\int h-g\,d\mu\le\epsilon$. \Hint{Ambil $\Phi$
sebagai himpunan fungsi-fungsi terintegralkan yang mempunyai sifat ini,
dan tunjukkan bahwa himpunan tersebut memenuhi syarat-syarat 122Yb.}
%122Yd, 122R

\spheader 122Yf Ulangi 122Yb dan 122Ye untuk ukuran Lebesgue pada
$\BbbR^r$, dengan $r>1$.
%122Yd, 122Yb, 122Ye, 122R

\spheader 122Yg Misalkan $(X,\Sigma,\mu)$ suatu ruang ukur, dan andaikan
bahwa terdapat sekurang-kurangnya satu partisi $X$ menjadi tak hingga
banyak himpunan terukur tak-kosong. Misalkan $f$ suatu fungsi bernilai
real yang didefinisikan pada subhimpunan koterabaikan dari $X$, dan
$a\in\Bbb R$. Tunjukkan bahwa pernyataan-pernyataan berikut ekuivalen:

\quad (i) $f$ terintegralkan, dengan $\int f=a$;

\quad (ii) untuk setiap $\epsilon>0$ terdapat suatu partisi
$\sequencen{E_n}$ dari $X$ menjadi himpunan-himpunan terukur tak-kosong
sedemikian sehingga

\Centerline{$\sum_{n=0}^{\infty}|f(t_n)|\mu E_n<\infty$,
\quad$|a-\sum_{n=0}^{\infty}f(t_n)\mu E_n|\le\epsilon$}

\noindent setiap kali $\sequencen{t_n}$ suatu barisan sedemikian sehingga
$t_n\in E_n\cap\dom f$ untuk setiap $n$. (Seperti biasa, tetapkan
$0\cdot\infty=0$ dalam rumus-rumus ini.) \Hint{gunakan 122Yd.}
%122Yd, 122R

\spheader 122Yh Temukan suatu rumusan kembali (g) yang mencakup
kasus ruang-ruang ukur yang {\it tidak} dapat dipartisi menjadi barisan
himpunan terukur tak-kosong.
%122Yg, 122Yd, 122R

\spheader 122Yi Misalkan $X$ suatu himpunan, $\Sigma$ suatu
aljabar-$\sigma$ atas subhimpunan-subhimpunan $X$, dan
$\sequencen{\mu_n}$ suatu barisan ukuran dengan domain $\Sigma$.
Tetapkan $\mu E=\sum_{n=0}^{\infty}\mu_nE$ untuk $E\in\Sigma$.
(i) Tunjukkan bahwa $\mu$ adalah suatu ukuran. (ii) Tunjukkan bahwa
untuk sembarang fungsi bernilai real $f$ yang didefinisikan pada suatu
subhimpunan $X$, $f$ terintegralkan terhadap $\mu$ jika dan hanya jika
ia terintegralkan terhadap $\mu_n$ untuk setiap $n$ dan
$\sum_{n=0}^{\infty}\int|f|d\mu_n$ berhingga, serta bahwa dalam hal itu
$\int fd\mu=\sum_{n=0}^{\infty}\int fd\mu_n$.
%122Xi

\spheader 122Yj Misalkan $X$ suatu himpunan, $\Sigma$ suatu
aljabar-$\sigma$ atas subhimpunan-subhimpunan $X$, dan
$\langle\mu_i\rangle_{i\in I}$ suatu keluarga ukuran dengan domain
$\Sigma$. Tetapkan $\mu E=\sum_{i\in I}\mu_iE$ untuk $E\in\Sigma$.
(i) Tunjukkan bahwa $\mu$ adalah suatu ukuran. (ii) Tunjukkan bahwa
untuk sembarang fungsi $\Sigma$-terukur $f:X\to\Bbb R$, $f$
terintegralkan terhadap $\mu$ jika dan hanya jika ia terintegralkan
terhadap $\mu_i$ untuk setiap $i$ dan
$\sum_{i\in I}\int|f|d\mu_i$ berhingga.
%122Yi, 122Xi
}%end of exercises

\endnotes{
\Notesheader{122} Seperti dalam \S121, beberapa masalah teknis tambahan
timbul karena saya bersikeras mencoba mengintegralkan (i) fungsi-fungsi
yang tidak didefinisikan pada seluruh ruang ukur yang sedang ditinjau;
(ii) fungsi-fungsi yang, dalam arti formal, tidak terukur, tetapi hanya
terukur pada suatu himpunan koterabaikan. Tidak ada sesuatu pun dalam
bagian ini yang membenarkan salah satu dari kedua perluasan tersebut.
Namun, dalam bagian berikutnya kita akan meninjau limit barisan fungsi,
dan limit-limit ini tidak harus terdefinisi di setiap titik; contoh yang
limitnya tidak terdefinisi pada semua titik sama sekali tidak bersifat
patologis, melainkan justru sentral bagi penerapan-penerapan teori yang
paling penting.

Persoalan mengintegralkan fungsi yang tidak sepenuhnya terukur lebih
terbuka untuk diperdebatkan. Saya akan membahas hal ini lebih lanjut
setelah memperkenalkan secara formal ruang ukur `lengkap' dalam Bab 21.
Untuk saat ini, saya hanya akan mengatakan bahwa menurut saya upaya
menyusun formalisasi yang, sejauh wajar, mengintegralkan sebanyak mungkin
fungsi memang layak dilakukan; salah satu tujuan semula integral Lebesgue
adalah memungkinkan kita mengintegralkan lebih banyak fungsi daripada
pendahulu-pendahulunya.

Definisi `integrasi' di sini berjalan melalui tiga tahap yang dapat
dibedakan: (i) integrasi fungsi sederhana (122A-122G); (ii) integrasi
fungsi nonnegatif (122H-122L); (iii) integrasi fungsi bernilai real yang
umum (122M-122R). Saya menempuh setiap tahap secara perlahan; misalnya,
saya baru beralih ke fungsi nonnegatif yang terintegralkan setelah
mempunyai perangkat lengkap lema yang diperlukan mengenai fungsi
sederhana. Tentu saja ada banyak sekali rute alternatif; lihat,
misalnya, 122Yd, yang menawarkan definisi dengan dua langkah alih-alih
tiga. Saya memilih pendakian yang lebih panjang dan landai ini sebagian
karena (menurut pandangan saya) ia memberi gambaran gagasan yang lebih
jelas, dan sebagian karena ia bersesuaian dengan metode yang hampir
kanonis untuk membuktikan sifat-sifat fungsi terintegralkan: kita
membuktikannya terlebih dahulu untuk fungsi sederhana, kemudian untuk
fungsi nonnegatif yang terintegralkan, lalu untuk fungsi terintegralkan
yang umum. (Petunjuk yang saya berikan untuk 122Yb mengikuti filosofi
ini. Lihat juga 122Xc; tetapi saya tidak menyatakannya sebagai teorema
formal, karena urutan pembuktiannya yang tepat berbeda dari satu kasus
ke kasus lain, dan menurut saya sebaiknya hal itu diingat sebagai metode
pendekatan, bukan sebagai suatu hasil khusus untuk diterapkan.)

Anda berhak merasa bahwa bagian ini sangat abstrak dan hampir tidak
memberi gambaran tentang fungsi-fungsi favorit Anda mana yang kemungkinan
terintegralkan, apalagi tentang nilai integralnya. Saya berharap Bab 13
akan sedikit membantu dalam hal ini, meskipun harus saya katakan bahwa
untuk metode yang benar-benar berguna dalam menghitung integral, kita
harus menunggu Bab 22, 25, dan 26 dalam jilid berikutnya. Jika Anda ingin
mengetahui pusat sejati argumen-argumen bagian ini, saya sendiri akan
menempatkannya pada 122G, 122H, dan 122K. Intinya, gagasan-gagasan
122A-122F berlaku bagi kelas struktur $(X,\Sigma,\mu)$ yang jauh lebih
luas, karena hanya melibatkan operasi pada sejumlah berhingga anggota
$\Sigma$ sekaligus; barisan himpunan sama sekali tidak muncul. Kunci
yang memungkinkan segala sesuatu sesudahnya adalah 122G, yang bertumpu
pada 112Cf. Dan setelah 122G-122K, bagian selebihnya, meskipun sama sekali
tidak elementer, sesungguhnya hanyalah serangkaian pemeriksaan cermat
untuk memastikan bahwa fungsional yang didefinisikan dalam 122K
berperilaku seperti yang kita harapkan.

Banyak hasil dalam bagian ini (termasuk hasil kunci, 122G) akan digantikan
oleh hasil yang lebih kuat dalam bagian berikutnya. Namun, saya perlu
menyebut Lema 122Ja, yang dari waktu ke waktu akan muncul kembali sebagai
kriteria yang sangat berguna bagi keterintegralan fungsi nonnegatif
(lihat 122Ra).

Ada satu hal lain mengenai integral standar sebagaimana didefinisikan di
sini. Integral ini merupakan integral `absolut', dalam arti bahwa jika
$f$ terintegralkan, maka $|f|$ juga terintegralkan (122P). Ini berarti
bahwa walaupun integral Lebesgue memperluas integral Riemann `wajar'
(lihat 134K di bawah), terdapat fungsi-fungsi dengan integral Riemann
`tak wajar' berhingga yang tidak terintegralkan secara Lebesgue; contoh lazimnya
ialah $f(x)=\bover{\sin x}x$, dengan
$\lim_{a\to\infty}\int_0^af$ ada di $\Bbb R$, sedangkan
$\lim_{a\to\infty}\int_0^a|f|=\infty$, sehingga $f$ tidak terintegralkan,
dalam arti yang didefinisikan di sini, pada seluruh interval
$\ooint{0,\infty}$. (Untuk pembuktian lengkap atas pernyataan-pernyataan
ini, lihat 283D dan 282Xm dalam Jilid 2.) Jika Anda pernah menjumpai teori
deret yang dapat dijumlahkan secara `absolut' dan `bersyarat', Anda tentu
mengetahui bahwa jenis yang kedua dapat memperlihatkan perilaku yang
sangat membingungkan, dan akan memahami bahwa membatasi pengertian
`terintegralkan' sehingga berarti `terintegralkan secara absolut' sangat
memudahkan.

Sesungguhnya, ini lebih dari sekadar kemudahan; pembatasan tersebut perlu
agar definisi dapat bekerja pada tingkat abstraksi yang digunakan dalam
bab ini. Tinjau contoh ukuran cacah $\mu$ pada $\Bbb N$ (112Bd, 122Xd).
Struktur $(\Bbb N,\Cal P\Bbb N,\mu)$ invarian terhadap permutasi; yakni,
$\mu(\pi[A])=\mu A$ untuk setiap $A\subseteq\Bbb N$ dan setiap permutasi
$\pi:\Bbb N\to\Bbb N$. Akibatnya, {\it setiap} definisi integrasi yang
hanya bergantung pada struktur $(\Bbb N,\Cal P\Bbb N,\mu)$ juga harus
invarian terhadap permutasi, yakni,

\Centerline{$\int f(\pi(n))\mu(dn)=\int f(n)\mu(dn)$}

\noindent untuk setiap fungsi terintegralkan $f$ dan setiap permutasi
$\pi$. Tetapi tentu saja (sebagaimana saya harap telah diberitahukan
kepada Anda) suatu deret $\sequencen{f(n)}$ sedemikian sehingga
$\sum_{n=0}^{\infty}f(\pi(n))=\sum_{n=0}^{\infty}f(n)\in\Bbb R$ untuk
setiap permutasi $\pi$ harus dapat dijumlahkan secara absolut. Jadi, jika
kita hendak mendefinisikan integral pada ruang ukur abstrak
$(X,\Sigma,\mu)$ dengan cara yang hanya bergantung pada $\Sigma$ dan
$\mu$, kita hampir tak terelakkan dipaksa mendefinisikan integral
absolut.

Tentu saja ada konteks-konteks penting tempat pembatasan ini mengganggu,
dan tempat suatu bentuk integral `tak wajar' tampak tepat. Contoh lazimnya
adalah teori transformasi Fourier, tempat kita meninjau
$\lim_{a\to\infty}\int_{-a}^af$ sebagai pengganti
$\int_{-\infty}^{\infty}f$ (lihat \S283). Sangat banyak bentuk integral
tak wajar yang lebih atau kurang abstrak telah diajukan; banyak yang menarik
dan beberapa penting; tetapi tidak satu pun menandingi integral `standar'
sebagaimana diuraikan dalam bab ini. (Untuk suatu upaya pemeriksaan
sistematis atas kelas tertentu integral tak wajar semacam itu, lihat Bab
48 dalam Jilid 4.)

Jauh lebih sedikit kajian yang telah dilakukan tentang integrasi
fungsi tak-terukur -- atau, lebih tepatnya, fungsi yang tidak sama hampir
di mana-mana dengan suatu fungsi terukur yang terintegralkan. Saya yakin
hal ini semata-mata karena terlalu sedikit masalah penting yang dapat
menunjukkan arah yang harus ditempuh. Dalam 134C di bawah, saya menyebut
pertanyaan apakah ada {\it suatu} fungsi bernilai real
tak-terukur pada $\Bbb R$. Jawaban standarnya adalah `ya', tetapi fungsi
semacam itu mustahil muncul sebagai hasil konstruksi biasa. Akibatnya,
sebagian besar pertanyaan mengenai fungsi tak-terukur muncul dalam
konteks-konteks yang sangat khusus, dan sejauh ini saya belum melihat
satu pun yang memberi petunjuk berguna tentang seperti apa perluasan yang
secara umum sesuai bagi pengertian `keterintegralan'.
}%end of notes

\discrpage
